\documentclass[a4paper]{article}
\usepackage{graphicx}
\usepackage{amsmath,amssymb}
\usepackage{minted}
\usepackage{multirow}
\usepackage{float}
\usepackage{algorithm}
\usepackage{algpseudocode}
\usepackage{todonotes}
\usepackage{forest}
\usepackage{xstring}
\usepackage{xcolor}
\usepackage[most]{tcolorbox}
\usepackage{xparse}
\usepackage{natbib}
\usetikzlibrary{graphs,graphdrawing}
\usegdlibrary{layered}
% ND Mind Mirror default input files.
\input{/home/donkarlo/Dropbox/repo/nd_language_ai_project/src/nd_language_ai/formal/computer/programming/tex/latex/code_clips/external_graphviz.tex}
\input{/home/donkarlo/Dropbox/repo/nd_language_ai_project/src/nd_language_ai/formal/computer/programming/tex/latex/part/control_sequence/kind/semtag/semtag.tex}
\input{/home/donkarlo/Dropbox/repo/nd_language_ai_project/src/nd_language_ai/formal/computer/programming/tex/latex/code_clips/seehere.tex}

\title{Mann--Kendall Trend Test}
\date{}
\author{Mohammad Rahmani}
\begin{document}
\maketitle

\section{Definition}
The Mann--Kendall test is a nonparametric test for evidence of a monotonic increasing or decreasing trend in an ordered sequence. It uses the signs of pairwise differences rather than requiring a linear relationship.

\section{Methodology}
For observations $y_1,\ldots,y_n$, define
\[
 S=\sum_{i=1}^{n-1}\sum_{j=i+1}^{n}
 \operatorname{sgn}(y_j-y_i)
\]
where $\operatorname{sgn}(u)$ is $1$ if $u>0$, $-1$ if $u<0$, and $0$ otherwise. Positive $S$ indicates more upward comparisons than downward comparisons.

With no tied values, under the no-trend null,
\[
 \operatorname{Var}(S)=\frac{n(n-1)(2n+5)}{18}.
\]
If the data contain tied groups with respective sizes $q_r$, use the tie correction
\[
 \operatorname{Var}(S)=\frac{n(n-1)(2n+5)
 -\sum_rq_r(q_r-1)(2q_r+5)}{18}.
\]
For a sufficiently large sample with positive variance, the continuity-corrected standardized score is
\[
 Z=\begin{cases}
 \dfrac{S-1}{\sqrt{\operatorname{Var}(S)}}, & S>0,\\
 0, & S=0,\\
 \dfrac{S+1}{\sqrt{\operatorname{Var}(S)}}, & S<0.
 \end{cases}
\]
A two-sided test evaluates evidence for a monotonic trend. A test statistic or $p$-value alone does not report the size of a trend.

\section{E-Commerce Example}
Apply the test to monthly sales \emph{after addressing strong seasonal patterns}. If evidence favors an upward monotonic trend, report its practical magnitude separately using a slope estimate such as Theil--Sen.

\section{Advantages}
\begin{itemize}
 \item Does not require normally distributed data.
 \item Less influenced by the magnitude of extreme values than OLS.
 \item Can detect monotonic but nonlinear changes.
\end{itemize}

\section{Limitations}
\begin{itemize}
 \item Serial correlation affects the null distribution; positive correlation can lead to false positives with naive inference.
 \item Strong seasonal cycles can obscure or mimic overall direction.
 \item Detects evidence of trend, not its slope or causal mechanism.
\end{itemize}
\section{Practical Guidance}
Check autocorrelation and seasonality before selecting the test. Consider seasonal Mann--Kendall or suitably adjusted inference for dependent observations. Report the period analyzed, $p$-value, and an interpretable estimated slope together.

\bibliographystyle{plainnat}
\bibliography{/home/donkarlo/Dropbox/repo/nd_shared_research/bibliography/bibliography.bib}
\end{document}
